\chapter{BCs and Loads}
\section{Boundary Conditions}
Boundary Conditions are used to describe the behavior of the system at
the boundary. Currently we support different types of boundary
conditions: 

\begin{itemize}
\item homogenous,
\item inhomogeneous and
\item absorbing BCs.
\end{itemize}

{\bf Homogenous BCs} are used to suppress the nodal DOF at the specified
points. Therewith, the value of the dof is set to zero. These
conditions are used, e.g. for grounding an electrode (electric potential
= 0).

{\bf Inhomogeneous BCs}  are used for assigning a
  specific value to a boundary. They describe non-zero (maybe
  time-varying) values.  Therefore, you have to select the name of the
  specific nodes (according to the section \xel{nodes} in
\xel{$<$domain$>$}) and a value which has to be assigned to this node.
  Additionally, a name of a time function file may be provided for
  time varying BCs. In this case, the resulting values applied to the
  selected nodes will be ${dof}(timestep) =
  function(timeFncName,timestep)*factor$, where $factor$ stands for
  the value given to the BCs. A typical application for an
  inhomogeneous BC is the excitation of an electric device by an
  electric potential.

{\bf Absorbing BCs} are used to get rid of the
  artificial reflection of waves at the simulation boundaries. For
  details see section "Acoustics".

\bigskip

In the section  \xel{bcsAndLoads} of the appropriate PDE in the
config-file (there can be more than one!), the BCs have to be defined
in the following manner
\begin{Verbatim}[commandchars=@\[\]]
  @PYaZ[<bcsAndLoads]@PYaZ[>]
     @PYaZ[<dirichletHom]   @PYaQ[name=]@PYaB["nameBound1"]@PYaZ[/>]
     @PYaZ[<dirichletHom]   @PYaQ[name=]@PYaB["nameBound2"]@PYaZ[/>]
     @PYaZ[<dirichletInHom] @PYaQ[name=]@PYaB["nameBound3"] @PYaQ[value=]@PYaB["-1.0"]@PYaZ[/>]
  @PYaZ[</bcsAndLoads>]
\end{Verbatim}


With inhomogeneous BCs, one has additionally to define the \xatv{value} by
giving the name of the boundary followed by the boundary value. If a
time varying signal has to be applied, one has additionally to provide
the name of the time file in the attribute {\xatv{signal.dat} }:
\begin{Verbatim}[commandchars=@\[\]]
  @PYaZ[<dirichletInHom] @PYaQ[name=]@PYaB["nameBound3"] @PYaQ[value=]@PYaB["10 * sample1D('signal.dat',t,1)"]@PYaZ[/>]
\end{Verbatim}


\noindent
If we want to define a boundary condition for a PDE with more
than one dof, we have additionally to define the dof, the BC has to
work on. For the homogenous BC one has to use
\begin{Verbatim}[commandchars=@\[\]]
  @PYaZ[<dirichletHom]   @PYaQ[name=]@PYaB["nameBound1"] @PYaQ[dof=]@PYaB["XXX"]@PYaZ[/>]
\end{Verbatim}

\noindent
for the inhomogenous BC the command

\begin{Verbatim}[commandchars=@\[\]]
  @PYaZ[<dirichletInHom]   @PYaQ[name=]@PYaB["nameBound1"]  @PYaQ[dof=]@PYaB["XXX"]@PYaZ[/>]
\end{Verbatim}


\noindent
has to be called. Herein, \xatv{XXX} stands for the name degrees of
freedom, e.g. in the mechanical PDE

\begin{Verbatim}[commandchars=@\[\]]
@PYaZ[<bcsAndLoads]@PYaZ[>]
   @PYaZ[<dirichletHom]   @PYaQ[name=]@PYaB["leftBound"]   @PYaQ[dof=]@PYaB["x"]@PYaZ[/>]
   @PYaZ[<dirichletHom]   @PYaQ[name=]@PYaB["leftBound"]   @PYaQ[dof=]@PYaB["y"]@PYaZ[/>]
   @PYaZ[<dirichletInHom] @PYaQ[name=]@PYaB["rightBound"]  @PYaQ[dof=]@PYaB["y"]  @PYaQ[value=]@PYaB["2.0"]@PYaZ[/>]
   @PYaZ[<load]           @PYaQ[name=]@PYaB["centerPoint"] @PYaQ[dof=]@PYaB["x"]  @PYaQ[value=]@PYaB["3.0"]@PYaZ[/>]
@PYaZ[</bcsAndLoads>]
\end{Verbatim}


\section{Loads}

{\bf Loads} are used to apply external nodal loads like an electric
charge.  The syntax of the {\xel{load} } command is equal to that of 
the {\em inhomogeneous BCs}:
\begin{Verbatim}[commandchars=@\[\]]
@PYaZ[<bcsAndLoads]@PYaZ[>]
   @PYaZ[<load] @PYaQ[name=]@PYaB["loadBound"] @PYaQ[value=]@PYaB["10"] @PYaQ[dynamics=]@PYaB["timeFile"]@PYaZ[/>]
@PYaZ[</bcsAndLoads>]
\end{Verbatim}


Similar to the boundary conditions, the degree of freedom (dof) to
which the load has to be applied, must be defined for PDEs with more
than one dof per node:

\begin{Verbatim}[commandchars=@\[\]]
  @PYaZ[<load]  @PYaQ[name=]@PYaB["nameBound1"]  @PYaQ[dof=]@PYaB["XXX"]@PYaZ[/>]
\end{Verbatim}





